% =====================================================================
%  Appendix
% =====================================================================
\appendix

% =====================================================================
\section{Proof of Theorem~\ref{thm:stability}}
\label{app:proof-stability}

This appendix proves the per-epoch and cumulative bounds of
Theorem~\ref{thm:stability}, together with the score-transfer
corollary used in Sec.~\ref{sec:method-core}. The proof has three
steps: a Davis--Kahan angular bound, a conversion to Euclidean
distance via sign calibration, and a transfer from anchor stability
to the \textsc{TADS} score.

\paragraph{Assumption interpretations.}
Assumption~1 holds when parameter updates are learning-rate bounded
and the map $\theta \mapsto h_l^{\mathrm{last}}(\cdot;\theta)$ is
locally Lipschitz. Re-sampling the probe can add an
$\eta_t$-independent term of order
$O(|\widetilde{\mathcal{D}}|^{-1/2})$; empirically this term is small
at the probe sizes used in our verification. Assumption~2
($\gamma^{(t)}\ge\gamma_{\min}$) is checked in
App.~\ref{app:thm-verification}. Assumption~3 is enforced by sign
calibration in Alg.~\ref{alg:tads}.

\subsection{Step 1: Angular bound}
\label{app:step-angular}

The Davis--Kahan $\sin\Theta$ theorem~\cite{daviskahan} states that
for symmetric matrices $A,B$ with top unit eigenvectors $u,v$ and
top eigengap $\delta=\lambda_1(A)-\lambda_2(A)>0$,
\[
\sin\Theta(u,v)\le \frac{\|A-B\|_F}{\delta},
\qquad
\Theta(u,v)=\arccos|\langle u,v\rangle|.
\]
Apply this with $A=\Sigma^{(t)}$, $B=\Sigma^{(t+1)}$,
$\delta=\gamma^{(t)}\ge\gamma_{\min}$, and
$\|\Sigma^{(t+1)}-\Sigma^{(t)}\|_F\le C_\Sigma\eta_t$. Then
\begin{equation}
  \sin\Theta\bigl(v_l^{(t+1)},v_l^{(t)}\bigr)
  \le
  \frac{C_\Sigma}{\gamma_{\min}}\eta_t .
  \label{eq:angular-bound}
\end{equation}

\subsection{Step 2: Per-epoch and cumulative bounds}
\label{app:step-euclidean}

The angular bound is invariant to the PCA sign ambiguity. Sign
calibration ensures
$\langle v_l^{(t+1)},v_l^{(t)}\rangle>0$, so
$\Theta\in[0,\pi/2)$. Hence
\begin{equation}
  \bigl\|v_l^{(t+1)}-v_l^{(t)}\bigr\|
  =
  2\sin(\Theta/2)
  \le
  2\sin\Theta
  \le
  \frac{2C_\Sigma}{\gamma_{\min}}\eta_t .
  \label{eq:perepoch-bound}
\end{equation}
This is the per-epoch bound in Theorem~\ref{thm:stability}. Summing
over epochs gives
\begin{equation}
  \bigl\|v_l^{(t+m)}-v_l^{(t)}\bigr\|
  \le
  \frac{2C_\Sigma}{\gamma_{\min}}
  \sum_{s=t}^{t+m-1}\eta_s .
  \label{eq:cumulative-bound-app}
\end{equation}
Under a decaying learning-rate schedule, the tail
$\sum_{s\ge t}\eta_s$ shrinks with $t$, yielding the stated
$\varepsilon$-stationarity.

\subsection{Step 3: Score stability transfer}
\label{app:score-stability}

Write the \textsc{TADS} score as
\[
s_i^{(t)} = B_i^{(t)}A_i^{(t)},
\qquad
B_i^{(t)} := \mathcal{R}_i^{(t)}a_i^{(t)},
\qquad
A_i^{(t)} := 1+\lambda\widetilde{\mathrm{align}}_i^{(t)} .
\]
Here $B_i^{(t)}$ is the reward--advantage consensus base ranking and
$A_i^{(t)}\in[1,1+\lambda]$ is the anchor factor.

\begin{corollary}[Score stability transfer]
\label{cor:score-stability}
Assume $|B_i^{(t)}|\le B_{\max}$, the sequence-mean activations
satisfy $\|\bar h_l(x_i;\theta_{t-1})\|\le H_{\max}$, and the
min--max normalisation in Eq.~\ref{eq:align} has Lipschitz constant
$L_{\mathrm{norm}}$. Then the anchor-induced score change satisfies
\begin{equation}
  \bigl|
  B_i^{(t+1)}
  \bigl(A_i^{(t+1)}-A_i^{(t)}\bigr)
  \bigr|
  \le
  B_{\max}\lambda L_{\mathrm{norm}}H_{\max}
  \frac{2C_\Sigma}{\gamma_{\min}}\eta_t
  =
  O(\lambda\eta_t).
  \label{eq:cor-score-bound}
\end{equation}
\end{corollary}

\begin{proof}
By Cauchy--Schwarz,
$\langle \bar h_l(x_i;\theta_{t-1}),v_l^{(t)}\rangle$ is
$H_{\max}$-Lipschitz in $v_l^{(t)}$. Therefore the raw alignment
changes by at most
$H_{\max}\|v_l^{(t+1)}-v_l^{(t)}\|$. Composing with the
$L_{\mathrm{norm}}$-Lipschitz min--max normalisation and applying
Eq.~\ref{eq:perepoch-bound} gives
\[
|A_i^{(t+1)}-A_i^{(t)}|
\le
\lambda L_{\mathrm{norm}}H_{\max}
\frac{2C_\Sigma}{\gamma_{\min}}\eta_t .
\]
Multiplying by $|B_i^{(t+1)}|\le B_{\max}$ proves
Eq.~\ref{eq:cor-score-bound}.
\end{proof}

The corollary holds for any bounded base ranking score
$B_i^{(t)}$. Thus the stability result concerns the trajectory
anchor and its multiplicative factor; it is not tied to a particular
reward, advantage estimator, or reward--advantage consensus rule.

% =====================================================================
\section{Experimental Setup Details}
\label{app:setup-details}

\paragraph{Hardware.}
Experiments run on NVIDIA A-series 80GB GPUs with NVLink. The main
LLaMA-2-7B experiments use a single multi-GPU node; the Qwen2.5-14B
LoRA scaling study uses an $8\times$A100 node.

\paragraph{Main experiment: LLaMA-2-7B full fine-tuning.}
The main experiments fine-tune all parameters of
LLaMA-2-7B~\cite{touvron2023llama2}. Training uses Adam
($\beta_1=0.9$, $\beta_2=0.999$), learning rate $2\times10^{-5}$,
cosine decay to zero, warmup ratio $0.03$, weight decay $0.1$, bf16
precision, 3 epochs, batch size 128, and cutoff length 2{,}048. The
configuration is held fixed across selection methods.

\paragraph{Evaluation protocol.}
We evaluate with lm-evaluation-harness~\cite{eval-harness} on nine
tasks: MMLU, MMLU-Pro, GSM8K, SVAMP, HumanEval, MBPP, TydiQA,
XQuAD, and BBH. Scores are reported per task and averaged into AVG.

\paragraph{\textsc{TADS}-specific hyperparameters.}
The default probe size is
$|\widetilde{\mathcal{D}}_t|=1{,}024$. The trajectory anchor is
refreshed each epoch, alignment is averaged over all transformer
layers, and the default anchor weight is $\lambda=1.0$.

\paragraph{Reward--advantage consensus base.}
At each refresh step $t$ we form a composite reward
\[
\mathcal{R}_i^{(t)}
=
w^{(t)}\mathcal{L}_i^{(t)}
+
\bigl(1-w^{(t)}\bigr)\mathcal{H}_i^{(t)},
\]
with $w^{(t)}$ defined in Eq.~\ref{eq:w}, and a per-sample
\emph{advantage estimate} $a_i^{(t)} \in (0,1)$.

\paragraph{Advantage-estimator architecture and training
(implementation device).}
We obtain $a_i^{(t)}$ from a small actor--critic network: a
two-layer MLP trunk ($256{\to}256$, ReLU) feeds a
$\mathrm{Beta}(\alpha_t(h_i),\beta_t(h_i))$ actor head and a scalar
critic head; the actor's bounded sample
$a_i \sim \mathrm{Beta}(\cdot,\cdot) \in (0,1)$ is our advantage
estimate. The network is trained with a clipped PPO surrogate on
$\mathcal{R}_i^{(t)}$. PPO here is an implementation device for
fitting a bounded per-sample advantage estimator -- any procedure
that produces a learned, non-negative, bounded $a_i$ correlated
with $\mathcal{R}_i$ would substitute. The selection rule
\eqref{eq:score} consumes $a_i$ only as a scalar value.

\begin{table}[h]
\centering
\small
\begin{tabular}{@{}ll@{}}
\toprule
hyperparameter & value \\
\midrule
trunk            & MLP $256 \to 256$, ReLU \\
heads            & $\alpha,\beta$ (softplus${+}1$), $V$ (linear) \\
actor distribution & $\mathrm{Beta}(\alpha,\beta)$, sample clamped to $[10^{-6}, 1{-}10^{-6}]$ \\
advantage mode   & group-relative (default); GAE available \\
GAE $\gamma$     & $0.99$ \\
GAE $\lambda$    & $0.95$ \\
PPO clip $\epsilon$ & $0.2$ (actor and value) \\
PPO epochs       & $4$ \\
mini-batch       & $1024$ \\
learning rate    & $3\times 10^{-4}$ \\
entropy coef.    & $0.01$ \\
value coef.      & $0.5$ \\
grad-norm clip   & $1.0$ \\
\bottomrule
\end{tabular}
\end{table}

\noindent Candidates are ranked by the reward--advantage consensus
score $\mathcal{R}_i^{(t)} a_i^{(t)}$ at $\lambda=0$, and by the
full \textsc{TADS} score \eqref{eq:score} at $\lambda>0$.

% =====================================================================
\section{$\lambda$-Ablation and Qwen2.5-0.5B Sanity Track}
\label{app:setup-validation}

Before the main experiments, we validate the pipeline on
Qwen2.5-0.5B. The setup uses full fine-tuning on Alpaca-GPT4
($52{,}002$ pairs), selection ratio $\rho=10\%$, 3 epochs, AdamW,
learning rate $2\times10^{-5}$, cosine schedule with $10\%$ warmup,
and evaluation on MMLU and GSM8K.

The $\lambda$-ablation is reported in Table~\ref{tab:lambda-05b} in
the main text. The curve is non-monotonic: the no-anchor consensus
base is improved by small non-zero $\lambda$, while overly large
$\lambda$ values reduce the gain. This supports the interpretation
of the anchor as a bounded structural reweighting rather than a
replacement for the base ranking.

% =====================================================================
\section{Scaling Study: Qwen2.5-14B}
\label{app:scaling}

We additionally test \textsc{TADS} on Qwen2.5-14B with LoRA to check
whether the anchor remains inexpensive at larger scale.

\paragraph{Setup.}
Qwen2.5-14B is trained with LoRA ($r=16$, $\alpha=32$, dropout
$0.05$; target modules
$\{q,k,v,o,\mathrm{gate},\mathrm{up},\mathrm{down}\}$proj). The
candidate pool is Alpaca-GPT4 ($N=52{,}002$), with 3 epochs, batch
size 64, learning rate $2\times10^{-4}$, cosine schedule, 3\%
warmup, and $8\times$A100 GPUs. The compared settings are Random,
the no-anchor consensus base, \textsc{TADS}, and full-data
fine-tuning.

\paragraph{Cost.}
The per-epoch cost breakdown is reported in
Table~\ref{tab:cost-14b-main}. The forward pass and fine-tuning
update are shared with the no-anchor consensus base; the
\textsc{TADS}-only operations are trajectory-anchor PCA and
alignment scoring over cached hidden states. These add roughly
15 seconds to a $\sim$35-minute epoch, corresponding to less than
$0.5\%$ overhead.

\paragraph{Accuracy.}
The Qwen2.5-14B accuracy table will report Random, no-anchor
consensus base, \textsc{TADS}, and full-data fine-tuning under the
same evaluation suite as the main experiments.

% =====================================================================
\section{Empirical Verification of Theorem~\ref{thm:stability}}
\label{app:thm-verification}

This appendix verifies that the anchor-stability theorem describes
real training dynamics. We instrument a Qwen2.5-0.5B LoRA run on
Alpaca-GPT4 and record the quantities appearing in
Theorem~\ref{thm:stability}: eigengap, covariance drift, anchor
displacement, sign consistency, and learning rate.

\paragraph{Setup.}
The verifier runs every 25 optimiser steps on a fixed
2{,}000-sample probe. For each of the 24 decoder layers, it computes
the top two eigenvalues of the contextualisation-delta covariance,
the eigengap
$\gamma_l^{(t)}=\lambda_1^{(t)}-\lambda_2^{(t)}$, the covariance
drift $\|\Sigma_l^{(t+1)}-\Sigma_l^{(t)}\|_F$, and the anchor
displacement
$d_l^{(t)}=\|v_l^{(t+1)}-v_l^{(t)}\|$. The fixed probe removes
probe-resampling noise from the diagnostic.

\paragraph{Assumptions.}
The eigengap stays positive on the monitored layers, with effective
$\gamma_{\min}=0.050$. The trimmed covariance-drift estimate gives
$\hat C_\Sigma=5.78\times10^3$. Sign calibration is stable: the
measured sign-flip rate is $3.6\%$, below the $5\%$ tolerance used
for the probabilistic sign-calibration check.

\paragraph{Per-step bound.}
At the constant operating point $\eta=2\times10^{-4}$, the empirical
bound is
\[
\frac{2\hat C_\Sigma}{\gamma_{\min}}\eta
=
\frac{2\cdot 5.78\times10^3}{0.050}\cdot 2\times10^{-4}
=
46.24 .
\]
Across $1{,}496$ layer-refresh measurements, the measured anchor
displacement never violates this bound. The median tightness is
$1.46\times10^{-3}$ and the worst-case tightness is
$7.4\times10^{-3}$.

\paragraph{No-anchor diagnostic.}
As a post-hoc comparison, we compute the same displacement
diagnostic on a matched no-anchor consensus control
(\texttt{use\_anchor=false}, $\lambda=0$). This control does not use
the anchor for selection, but its hidden-state trajectory can still
be measured. Compared with \textsc{TADS}, the no-anchor control has
larger and more variable displacement statistics and fails the
probabilistic sign-stability check, as summarised in
Table~\ref{tab:dstep-spread}.

% =====================================================================
\section{In-Progress Empirical Study}
\label{app:inprogress}

The broader empirical study extends the main experiments along three
axes. First, we scale across model sizes from Qwen2.5-0.5B to the
7B-class main setting and Qwen2.5-14B. Second, we test robustness
across additional instruction-tuning corpora beyond Alpaca-GPT4,
including Evol-Instruct--style data. Third, we run additional
ablations over selection ratio $\rho$, anchor refresh frequency,
layer subsets, sign calibration, and probe size.

The goal is to test whether a single \textsc{TADS} configuration
continues to provide gains across model scales, datasets, and
benchmark families, and whether the trajectory anchor remains a
useful signal when the candidate pool distribution changes.
